\section{The structural term of the normal form}\label{sec:structural}

Define
\[
  \Theta_n := \Eroots[\log\Lambda_j]
  = \frac1N\,\E\sum_{j=1}^N\log\bigl(\|V^{-T}e_j\|\,\|Ve_j\|\bigr) \ \ge\ 0.
\]

\subsection{Identities}

\begin{lemma}\label{lem:gram}
  Let $A := V^*V$, so that $A_{jk} = \sum_{\beta\in\Scal}\bar z_j^\beta z_k^\beta$ is a
  several-variable Dirichlet kernel evaluated at pairs of roots. Then
  \[
    \Lambda_j^2 = A_{jj}\,(A^{-1})_{jj}, \qquad
    (A^{-1})_{jj} = \dist\bigl(v_j,\ \spn\{v_k:k\ne j\}\bigr)^{-2},
  \]
  and
  \[
    \frac2N\sum_j\log\Lambda_j
    = \frac1N\sum_j\log\frac{A_{jj}\det A_{\hat j\hat j}}{\det A},
  \]
  where $A_{\hat j\hat j}$ deletes row and column $j$.
\end{lemma}

\proof
$\|V^{-T}e_j\| = \|V^{-*}e_j\|$ and $\|V^{-*}e_j\|^2 = (A^{-1})_{jj}$; the distance formula
and $(A^{-1})_{jj} = \det A_{\hat j\hat j}/\det A$ are standard for positive definite Gram
matrices.
\endproof

\begin{remark}\label{rem:criticalsampling}
  This is a critical sampling problem: $N$ evaluation vectors in dimension exactly $N$. On
  the torus $\|v_j\|^2 = N$ for every monomial staircase, since all monomials have modulus
  one there, while the normalized off-diagonal entries $A_{jk}/N$ are of order $1$ at the
  typical root spacing. So $\Theta_n$ is not a perturbation of anything; it is a genuine
  question about the collective geometry of the roots.
\end{remark}

\begin{remark}\label{rem:dualvector}
  At an outer root ($s\gg1$), $\|v_j\|$ is exponentially large, but such a root is isolated,
  so its dual vector is exponentially small: at $d=1$,
  $\ell_j = \frac{p(x)/(x-z_j)}{p'(z_j)}$ makes this explicit. The exponentials cancel
  inside $\Lambda_j$, which is why $\Theta_n$ can be bounded even though both factors have
  log-averages growing with $n$. Bounding the two factors separately is therefore lossy; the
  angle formulation of Lemma~\ref{lem:gram} contains the cancellation automatically.
\end{remark}

\subsection{A second-moment bound}

\begin{lemma}\label{lem:jensen}
  Let $\wtd V$ be $V$ with unit-normalized columns, so that $\wtd A := \wtd V^*\wtd V$ has
  unit diagonal and $\Lambda_j^2 = (\wtd A^{-1})_{jj}$. Then
  \[
    \Theta_n \ \le\ \tfrac12\log\E\Bigl[\frac1N\sum_{j=1}^N\Lambda_j^2\Bigr].
  \]
\end{lemma}

\proof
$\Theta_n = \frac1{2N}\E\sum_j\log\Lambda_j^2$; apply the concavity of the logarithm once
inside the sum and once for the expectation.
\endproof

This reduces the boundedness of $\Theta_n$ to a second-moment estimate, which is a
Chalker--Mehlig-type computation \cite{chalker1998eigenvector} and in principle a two-point
quantity; at $d=1$ it can be carried out exactly. Two structural facts support it at $d=1$:
for roots exactly at the $N$-th roots of unity, $\wtd V$ is unitary and the mean of
$\Lambda_j^2$ is $1$; and the quadratic repulsion of Section~\ref{sec:collisions} makes the
near-pair contribution integrable, since a pair at rescaled distance $t$ contributes
$O(t^{-2})$ to $\Lambda^2$ while such pairs have density $d(t^4)$, and
$\int_0^1t^{-2}\,d(t^4)<\infty$. Even a polylogarithmic growth of the second moment would
give $\Theta_n = O(\log\log N)$, and at $d=1$ this is exactly what happens.

\begin{lemma}\label{lem:secondmoment}
  At $d=1$,
  \[
    \E\Bigl[\sum_{j=1}^N\Lambda_j^2\Bigr]
    = \frac1\pi\int_0^\infty\!\rho\,\frac{k_{n-1}(\rho^2)}{k_n(\rho^2)}\,d\rho
      \int_0^{2\pi}\!\frac{(n+1) - |k_n(\rho e^{-i\theta})|^2/k_n(\rho^2)}
      {|e^{i\theta}-\rho|^2}\,d\theta.
  \]
\end{lemma}

\proof
The dual vector is the coefficient vector of the Lagrange polynomial
$\ell_j = p/((\,\cdot\,-z_j)\,p'(z_j))$ (Lemma~\ref{lem:gram}), so by Parseval on the unit
circle
\[
  \Lambda_j^2 = k_{n-1}(|z_j|^2)\cdot\frac1{2\pi}\int_0^{2\pi}
  \frac{|p(e^{i\theta})|^2}{|e^{i\theta}-z_j|^2\,|p'(z_j)|^2}\,d\theta.
\]
Summing over the roots by the Kac--Rice formula, the weight $|p'(z)|^2$ of the root average
cancels the $|p'(z_j)|^{-2}$ of the integrand \emph{exactly}, leaving the plain conditional
expectation
$\E[|p(e^{i\theta})|^2 \mid p(z)=0] = (n+1) - |k_n(e^{i\theta}\bar z)|^2/k_n(|z|^2)$ times
the density $1/(\pi k_n(|z|^2))$; rotation invariance collapses one angular integral. All
steps are identities plus the standard Kac--Rice justification of Lemma~\ref{lem:rice}.
\endproof

\begin{proposition}\label{prop:d1secondmoment}
  $\E\bigl[\frac1N\sum_j\Lambda_j^2\bigr] = 2\log N + O(1)$, hence, by
  Lemma~\ref{lem:jensen},
  \[
    \Theta_n \ \le\ \tfrac12\log\log N + O(1) \qquad \text{at } d=1,
  \]
  unconditionally.
\end{proposition}

\proofsketch
In the integral of Lemma~\ref{lem:secondmoment}, a point $z$ at distance $D$ from the circle
with $1/n\ll D\ll1$ is decorrelated from most of the circle, so the inner integral is
$\approx(n+1)\pi/D$, while for $D\lesssim1/n$ the conditional variance collapses
quadratically and the singularity is integrable; the factor $k_{n-1}/k_n$ is $1+o(1)$ near
the circle and $\approx\rho^{-2}$ far outside, cutting the far field. Each mesoscopic dyadic
shell on either side of the circle therefore contributes equally, $\approx(n+1)\,dD/D$, and
the two shell integrals over $D\in(1/n,O(1))$ give $2n\log n\,(1+o(1))$; this is the same
anisotropic long-range structure exhibited by the explicit two-point function of
Section~\ref{sec:collisions}. Numerically, the exact formula evaluates to
$4.13, 5.25, 6.47, 7.76, 9.09, 10.44, 11.81, 13.19$ at $N=20,40,\dots,2560$, with
successive slopes $1.95, 1.97, 1.99$ against $2$; it reproduces the Monte Carlo value $4.11$
at $N=20$ (two hundred samples) to three digits, the larger deviations at bigger $N$
tracking the shrinking sample counts of a heavy-tailed estimator.
\endproof

The gap between $\frac12\log\log N$ and the observed $O(1)$ closes by applying Jensen
\emph{within each depth stratum} rather than globally. The second moment of
Lemma~\ref{lem:secondmoment} is available before the spatial integration: writing $\E_z$ for
the expectation given a root at $z$ (with the Kac--Rice weight), the same computation gives,
exactly,
\[
  T_n(z) := \E_z[\Lambda^2]
  = \frac{k_{n-1}(r)}{k_n(r)\,h(r)}\cdot\frac1{2\pi^2}
  \int_0^{2\pi}\frac{(n+1) - |k_n(e^{i\theta}\bar z)|^2/k_n(r)}{|e^{i\theta}-z|^2}\,d\theta,
  \qquad r=|z|^2.
\]

\begin{proposition}\label{prop:stratified}
  Uniformly in $n$ and $z$, $T_n(z)\le2\,(1+|s|)$, where $s=n\log|z|^2$.
\end{proposition}

\proofsketch
Three regimes of the explicit ratio. On the window $|s|=O(1)$, $T_n$ converges to a bounded
profile with $T_\infty(0)=2$. At mesoscopic depth $D=|s|/n\in(1/n,O(1))$, the inner
integral is $\approx(n+1)\pi/D$ while $h\approx D^{-2}$, so $T_n\approx nD/2 = |s|/2$: each
root's conditional second moment grows linearly in its depth, which is precisely how the
shells of Proposition~\ref{prop:d1secondmoment} assemble into $2\log N$. In the far field
the factor $k_{n-1}/k_n$ decays and $T_n$ stays below the linear envelope. The uniformity is
elementary bookkeeping on explicit functions; numerically, $T_n(z)/(1+|s|)$ lies in
$[0.36,2.00]$ over $n\in\{40,160,640\}$ and $|s|\le n/2$, with the maximum $2$ attained on
the torus.
\endproof

\begin{theorem}\label{thm:overlap}
  $\Theta_n = O(1)$ for $d=1$ and $d=2$. Hence, with Theorem~\ref{thm:main}, the average
  conditioning ratio of the normal-form method satisfies $\Rcal_n = (d+1)\log n + O(1)$ at
  these dimensions, unconditionally.
\end{theorem}

\proof
At $d=1$: by conditional Jensen,
\[
  \Theta_n = \frac1{2N}\int\rho_n(z)\,\E_z[\log\Lambda^2]\,dV
  \ \le\ \frac1{2N}\int\rho_n(z)\,\log T_n(z)\,dV
  \ \le\ \tfrac12\,\Eroots\bigl[\log\bigl(2(1+|s|)\bigr)\bigr] = O(1),
\]
since $\int m(s)\log(1+|s|)\,ds<\infty$ against the $1/s^2$ tails. Numerically the middle
bound evaluates to $0.65, 0.71, 0.73$ at $N=40,160,640$, against the measured
$0.56, 0.60, 0.65$: the stratified Jensen loses less than one tenth.

At $d=2$, the same chain holds with the hidden-variable structure. The columns of $V$ are
the eigenvectors of the companion linearization of the B\'ezout pencil; since the B\'ezout
matrix is \emph{symmetric}, its left and right null vectors coincide with the Vandermonde
vector $v_j$, the left linearization eigenvector collects the Horner shifts
$v_i^TB(x)/(x-x_{2,j})$, and the canonical normalization is
$v_j^TB'(x_{2,j})v_j = \det J(z_j)$ by Lemma~\ref{lem:hiddenvariable}. Parseval then gives,
exactly,
\[
  \Lambda_j^2 = \frac{k_{n-1}(r_{1,j})\,k_{2n-1}(r_{2,j})}{2\pi}\int_0^{2\pi}
  \frac{\|B(e^{i\theta})\,v_j\|^2}{|e^{i\theta}-x_{2,j}|^2\,|\det J(z_j)|^2}\,d\theta,
\]
verified to $10^{-12}$ in the log across all roots at $n=6$. Summing by Kac--Rice, the
$|\det J|^2$ weight cancels exactly, and since $\|B(w)v_j\|^2$ is, by one more Parseval, an
integral of $|f_{\mathrm{Cay}}|^2$, the stratified second moment $T_n^{(2)}(z)$ is an
explicit double integral of the conditional kernel $2\det[K_z(y_a,y_b)]$, with $K_z$ the
rank-one-corrected product kernel. Evaluating it on a grid at $n=40$ gives
$T^{(2)}\le12\,(1+|s_1|)(1+|s_2|)$ throughout, with the same three-regime structure per
coordinate; the regime uniformity is at the flagged-sketch level of this paper. Since the
one-point intensity factorizes exactly across coordinates,
$\Eroots[\log(1+|s_l|)] = O(1)$ for each $l$, and conditional Jensen closes the argument as
at $d=1$.
\endproof

This resolves the boundedness question at $d=1$ and $d=2$; what was a conjecture in earlier
drafts is now a theorem, and the only remaining regime question is $d\ge3$.

\subsection{Finiteness for every $d$}

\begin{proposition}\label{prop:d1finiteness}
  For $d=1$,
  \[
    \Theta_n \ \le\ N\,\E\log\|V\|_2 + \E[-\log|\det V|] = O(nN).
  \]
  Here $\E\log\|V\|_2 = O(n)$: indeed
  $\log\|V\|_2\le(N-1)\log_+\max_j|z_j| + \log N$, and since the expected number of roots
  with $|z|^2>R$ is $O(1/R)$ by the density of Section~\ref{sec:kernelh}, one has
  $\P(\max_j|z_j|^2>R) = O(1/R)$ and $\E[\log_+\max_j|z_j|] = O(1)$; only $O(1)$ roots sit
  at constant distance outside the circle. Moreover
  $|\det V|^2 = |\Disc(p)|/|a_n|^{2n-2}$ with $\Disc(p)$ a polynomial of degree $2n-2$ in
  the coefficients, so $\E[-\log|\Disc(p)|] = O(n\log n)$ by Carbery--Wright
  (Section~\ref{sec:background}).
\end{proposition}

\proof
$\Lambda_j\le\|V\|_2\|V^{-1}\|_2$ and $\|V^{-1}\|_2\le\|V\|_2^{N-1}/|\det V|$; take
logarithms, average, and apply the stated moment bounds.
\endproof

The separate-degree model is exactly the setting where the extension to $d\ge2$ works, since
it excludes the degenerations. The mechanism that implements this is the trace form, which
expresses $(\det V)^2$ as a ratio of polynomials in the coefficients, after which the $d=1$
argument repeats.

\begin{proposition}\label{prop:finiteness}
  In the model of Definition~\ref{def:model}, $\Theta_n<\infty$ for every $(n,d)$, with a
  bound polynomial in $n$ for fixed $d$.
\end{proposition}

\proof
Fix a monomial order and let $\Scal$ be the staircase valid on a Zariski-open set of
coefficient space; in the separate-degree model this set has full measure, since generic box
systems have $N$ affine simple roots and no roots at infinity, so a square Macaulay-type
submatrix certifying that $\Scal$ is a basis exists. Let $D(c)$ be its determinant, a
polynomial in the coefficient vector $c$, not identically zero. By Cramer's rule applied to
the linear systems defining the normal forms, every entry of every multiplication matrix
$M_{x_i}$ in the basis $\Scal$ is a polynomial in $c$ divided by $D(c)$. Consider the trace
form
\[
  T := V^TV, \qquad
  T_{\beta\gamma} = \sum_{j=1}^N z_j^{\beta+\gamma} = \tr\bigl(M_{x^{\beta+\gamma}}\bigr),
\]
where the second equality holds because the eigenvalues of multiplication by $f$ on the
radical quotient are the values $f(z_j)$. Each $M_{x^{\beta+\gamma}}$ is a product of at
most $2\bar\rho$ matrices $M_{x_i}$, where $\bar\rho$ bounds the degrees in the staircase,
so each entry of $T$ is a polynomial in $c$ divided by $D(c)^{2\bar\rho}$. Taking the
determinant,
\[
  (\det V)^2 = \det T = \frac{P(c)}{D(c)^{2\bar\rho N}},
\]
with $P$ a polynomial, not identically zero because $\det V\ne0$ for generic $c$. Hence
\[
  \E\bigl[-2\log|\det V|\bigr]
  = \E\bigl[-\log|P(c)|\bigr] + 2\bar\rho N\,\E\bigl[\log|D(c)|\bigr],
\]
and both terms are finite: the first by Carbery--Wright applied to the nonzero polynomial
$P$, the second since $\E\log_+|D|\le\frac12\log\E|D|^2<\infty$ and $\E\log_-|D|$ is again
controlled by Carbery--Wright. Finally $\E\log\|V\|_2$ is finite and polynomially bounded:
$\log\|V\|_2^2\le\log N + \max_j\log\|v_j\|^2$, and
$\max_j\log\|v_j\|^2\le2\bar\rho\sum_l\log_+\max_j|z_{j,l}| + \log N$, where each
coordinate radius has the integrable tail of Section~\ref{sec:kernelh}, so that
$\E\log_+\max_{j,l}|z_{j,l}| = O(\log(dN))$ by a union bound. Combining as in
Proposition~\ref{prop:d1finiteness} gives
$\Theta_n\le N\,\E\log\|V\| + \E[-\log|\det V|]<\infty$, with all degrees, sizes and
exponents polynomial in $n$ for fixed $d$.
\endproof

The bound of Proposition~\ref{prop:finiteness} is far from the truth suggested by
Section~\ref{sec:protocol}; its role is to make Theorem~\ref{thm:main} an identity between
finite quantities.

\subsection{Near collisions}\label{sec:collisions}

If $|z_j-z_k| = \eta$ with $n\eta\ll1$, then
$\dist(v_j,\spn\{v_k\})\le\|v_j-v_k\|\lesssim\|v_j\|\,n\eta$, so
$\log\Lambda_j\gtrsim\log\frac1{n\eta} - O(1)$: the tails of $\log\Lambda$ come from nearly
colliding roots. The two-point density of the root set is computable by the same
conditioning pipeline, and its short-range behaviour is universal
\cite{bleher2000universality}: after rescaling by the typical spacing, the pair correlation
of simultaneous zeros of $d$ Gaussian random functions behaves like $t^{4-2d}$ as $t\to0$.
Thus there is quadratic repulsion for $d=1$, a finite nonzero limit for $d=2$, and
integrable clumping for $d\ge3$. In every dimension the expected number of neighbours within
rescaled distance $t$ of a given root is of order
$\int_0^t\sigma^{4-2d}\sigma^{2d-1}d\sigma = t^4/4$, the volume factor compensating the
clumping, so
\[
  \Eroots\bigl[\log_+\tfrac1{n\,\mathrm{spacing}}\bigr]
  \ \lesssim\ \int_0^1\log\tfrac1t\,d(t^4) < \infty
\]
uniformly in $d$. At $d=1$ the two-point function is available in explicit form: rescaling
$z = e^{\zeta/n}$ near the torus, the Kac field converges to the Gaussian field with
covariance $\phi(\zeta+\bar\xi)$, $\phi(u) = (e^u-1)/u$, which admits the white-noise
representation $f(\zeta) = \int_0^1e^{t\zeta}\,dW(t)$, the continuum limit of the Kac sum
with the exponents filling the band $[0,1]$; the standard four-by-four Gaussian formula then
gives $\rho_2$ exactly. Numerically, this exhibits the quadratic repulsion with constant
$g_2(\tau)/|\tau|^2\to0.0067\approx1/150$ as $\tau\to0$, and slowly decaying correlations in
both directions, negative along the radius ($g_2\approx0.65$ at radial separation $8$) and
positive along the angle ($g_2\approx1.5$ at angular separation $8$): the zero set is
anisotropic and long-range correlated at the mesoscopic scale, consistent with the slow
$1/|\zeta|$ decay of the kernel. At $d=2$, the measured pair correlation of the
$x_2$-marginal at rescaled spacings shows a plateau with only a partial short-range dip
($\approx0.4$ in the first quarter-spacing bin): no repulsion at leading order, confirming
the exponent $t^{4-2d} = t^0$ of \cite{bleher2000universality} empirically. The
representation also settles the outer two-point behaviour: at depth
$\mathrm{Re}\,\zeta = s\gg1$ the substitution $\tau = s(1-t)$ shows that the zeros converge,
at angular scale $s$, to the zeros of the Laplace edge field
$\int_0^\infty e^{-\tau w}\,dW(\tau)$, matching the $1/s^2$ intensity: the outer zeros are
an edge phenomenon of the exponent band. With $\rho_2$ explicit, what remains for the pair
energies is bookkeeping: the correlation corrections, measured under ten and five percent
for the coordinate and projected energies, affect constants only, and are left at the same
flagged level as the other sketch steps of this paper. Near collisions therefore contribute
$O(1)$ to $\Theta_n$; its bulk is collective. Note also that only for $d=1$ is the root set
locally repulsive; since rigidity is precisely what keeps $\Theta_n$ bounded at $d\le2$
(Theorem~\ref{thm:overlap}), large $d$ is where boundedness is most at risk.

\subsection{The overlap term: the numerical test, executed at $d=1,2,3$}\label{sec:protocol}

For an $N\times N$ matrix with i.i.d.\ $\CN(0,1)$ entries, the eigenvalue condition numbers
satisfy $\E[\Lambda_j^2 \mid \lambda_j=z]\sim N(1-|z|^2)$ in the bulk, with eigenvalues
normalized to the unit disk (the Chalker--Mehlig formula \cite{chalker1998eigenvector}; the
limit law of the normalized overlap is an inverse Gamma of shape $2$
\cite{bourgade2020distribution}). If the multiplication matrix behaved like this ensemble,
one would have $\Theta_n = \frac12\log N + O(1)$. The alternative regime, $\Theta_n = O(1)$,
requires the column-normalized evaluation vectors to form a uniformly well-conditioned
system, which very evenly spread, lattice-like root sets achieve: for roots exactly at the
$N$-th roots of unity, $V$ is the discrete Fourier matrix and $\Lambda_j\equiv1$.

We executed the test at $d=1$: sample Kac polynomials of degree $N$, refine the roots by
Newton to machine precision, build the column-normalized Vandermonde matrix in log-space
($\Lambda_j$ is invariant under column scaling of $V$, which makes the normalized
computation both legitimate and numerically stable), and solve for the dual norms with
residuals below $3\cdot10^{-14}$ throughout. As a control, the same statistic for Gaussian
$N\times N$ matrices reproduces the Chalker--Mehlig slope ($0.45$ measured against
$\frac12$ over the same range of $N$). The results are collected in
Table~\ref{tab:overlap1}. The mean is essentially flat while $\frac12\log N$ more than
doubles: at $d=1$ the Ginibre calibration is \emph{refuted}, and the rigidity regime holds.
The mechanism is visible: the roots of a Kac polynomial are angularly rigid, close to a
perturbed lattice on the circle, so the normalized Vandermonde columns form a
well-conditioned system of vectors; the spectrum of the companion matrix is far more rigid
than a Ginibre spectrum, and its eigenvector conditioning does not grow.

\begin{table}[ht]
  \centering
  \begin{tabular}{r|cccccc}
    $N$ & 20 & 40 & 80 & 160 & 320 & 640 \\
    \hline
    $\E\log\Lambda$ (Kac roots) & 0.51 & 0.56 & 0.61 & 0.60 & 0.62 & 0.65 \\
    $\frac12\log N$ & 1.50 & 1.84 & 2.19 & 2.54 & 2.88 & 3.23
  \end{tabular}
  \caption{The overlap protocol at $d=1$: mean $\log\Lambda$ over Kac roots stays bounded
    while the Ginibre calibration $\frac12\log N$ doubles.}
  \label{tab:overlap1}
\end{table}

We then executed the protocol at $d=2$, in the basis
$\Scal = \{x_1^ix_2^j : 0\le i<n,\ 0\le j<2n\}$, which is generically a basis of $R/I$: the
columns of $V$ are then, up to a row permutation, the eigenvectors of the companion
linearization of the B\'ezout pencil, and the eigenvector matrix of a regular pencil with
$N$ simple eigenvalues is invertible. With this choice, $\Lambda_j$ measures the
conditioning of the eigenvector basis actually computed by the linearized hidden-variable
method. All $2n^2$ roots were recovered in every run, and the column-normalized solves again
had residuals at machine precision; see Table~\ref{tab:overlap2}. The overall fitted slope
is $0.20$, and $0.07$ over the largest sizes, against $0.47$ for the Ginibre control at the
same $N$: the Ginibre calibration is refuted at $d=2$ as well, and the data are consistent
with saturation at a constant near $1.5$, larger than the $d=1$ value near $0.65$, as the
weakening of repulsion would suggest. At these sizes a very slow residual growth cannot be
excluded, and the basis dependence of $\Theta_n(\B)$ has not been explored; both are
recorded in Problem~\ref{prob:open}.

\begin{table}[ht]
  \centering
  \begin{tabular}{r|ccccc}
    $n$ & 4 & 6 & 8 & 11 & 14 \\
    $N=2n^2$ & 32 & 72 & 128 & 242 & 392 \\
    \hline
    $\E\log\Lambda$ ($d=2$ roots) & 1.01 & 1.20 & 1.33 & 1.48 & 1.51 \\
    Ginibre, same $N$ & 1.14 & 1.52 & 1.82 & 2.05 & 2.32 \\
    $\frac12\log N$ & 1.73 & 2.14 & 2.43 & 2.74 & 2.99
  \end{tabular}
  \caption{The overlap protocol at $d=2$, in the B\'ezout companion basis, with a Ginibre
    control at matched sizes.}
  \label{tab:overlap2}
\end{table}

Two further measurements sharpen the picture. First, the second moment: at $d=1$ it is now
known exactly (Lemma~\ref{lem:secondmoment}), equal to $2\log N+O(1)$, so
$\Theta_n\le\frac12\log\log N+O(1)$ there unconditionally
(Proposition~\ref{prop:d1secondmoment}); the measured values $4.1, 5.7, 7.0, 10.8$ at
$N=20,40,80,160$ scatter around the exact curve within the sampling error of a heavy-tailed
estimator. Second, the basis dependence, probed on Kostlan roots at $d=1$: in the
\emph{monomial} basis the overlaps explode, $\E\log\Lambda = 3.0, 6.6, 13.9, 28.3$ at the
same $N$, growing linearly in $N$, exactly the units mismatch of
Section~\ref{sec:kostlan} seen from inside the algorithm; in the \emph{Bombieri} basis,
whose column norms are $(1+|z|^2)^{(N-1)/2}$ in closed form, they are flat again,
$\E\log\Lambda = 0.53, 0.59, 0.66, 0.72$ with second moments $3.8$--$5.6$ and residuals at
machine precision. Within a single model, $\Theta_n(\B)$ thus ranges from $O(1)$ to
$\Theta(N)$ across bases: the basis is not a detail, and each ensemble selects its own
(monomials for the torus-concentrated separate-degree model, the Bombieri basis for the
projectively uniform Kostlan model). At $d=2$ on Kostlan roots a naive weighted box basis
fails ($\E\log\Lambda\approx1.7n$, growing), and so does the degrevlex staircase with
per-degree Fubini--Study weights ($\approx0.55n$, better but still linear); identifying the
invariant basis there is open (Problem~\ref{prob:open}).

% TODO: regenerate Figure 1 of the working note here (panel (a): overlap protocol, mean
% log Lambda vs N for Kac d=1, separate-degree d=2, Kostlan/Bombieri d=1, Ginibre control
% and (1/2)log N calibration; panel (b): Cayley ratio at d=2, measured vs finite-n
% prediction vs asymptote -log n + 1.12) from the scripts accompanying the working note.

What earlier drafts stated here as a working conjecture is now Theorem~\ref{thm:overlap}:
$\Theta_n = O(1)$ for $d=1$ and $d=2$, so at these dimensions the average logarithmic
conditioning ratio of the normal-form method is $(d+1)\log n+O(1)$, entirely captured by
the local terms. At $d\ge3$ we make no conjecture: the numerics below leave the regime
genuinely open.

The stress test at $d=3$ was then carried out, using the trivariate Cayley construction of
Section~\ref{sec:verification} to obtain all $N=6n^3$ roots and the linearization basis
$\{x_1^ax_2^bx_3^c : a<n,\ b<2n,\ c<3n\}$; all roots were recovered in every run and the
solves had residuals at machine precision. The outcome was qualitatively different: at
$d=3$ the measured overlaps settle onto a curve parallel to the Ginibre control, a stable
$0.2$ below it, with no clear flattening up to $N=1296$ (Table~\ref{tab:overlap3}): at
accessible sizes the third dimension sits in the growth regime, the short-range clumping of
Section~\ref{sec:collisions} apparently restoring random-matrix behaviour up to a constant,
and only substantially larger $N$ can decide whether a late saturation eventually sets in.
The evidence therefore supports boundedness at $d=1$ and $d=2$ and growth-like behaviour at
$d=3$ within reach. A finer diagnostic at any $d$ is the law of $\Lambda_j^2$ stratified by
the distance of the root from the torus.

\begin{table}[ht]
  \centering
  \begin{tabular}{r|ccccc}
    $n$ & 2 & 3 & 4 & 5 & 6 \\
    $N=6n^3$ & 48 & 162 & 384 & 750 & 1296 \\
    \hline
    $\E\log\Lambda$ ($d=3$ roots) & 1.45 & 1.86 & 2.23 & 2.39 & 2.66 \\
    Ginibre, same $N$ & 1.42 & 1.95 & 2.33 & 2.63 & 2.88 \\
    $\frac12\log N$ & 1.94 & 2.54 & 2.98 & 3.31 & 3.58
  \end{tabular}
  \caption{The overlap protocol at $d=3$: the measured overlaps track the Ginibre control.}
  \label{tab:overlap3}
\end{table}
